% !PRINTED 248
\weberpage{248}
%<<CONT
reell, dagegen positiv, wenn $\alpha$ und $\gamma$ conjugirt (und dann wegen $\alpha+\gamma=0$ rein imaginär) sind. Daraus folgt:

% [corrected 2026-08-07, zoom adjudication pass] Both of these carried \gesperrt and should NOT:
% on the page the (8)/(9) pair is set in ordinary roman, tight, identical to the surrounding
% prose. The discriminating evidence is on this same page -- the grammatically parallel statement
% three lines below, before (10), genuinely IS spaced, "ist" included. So Weber spaces one and
% not the other, and the transcription applied Sperrung to (8)/(9) by analogy with (10)/(11)
% rather than by reading the type.
%
% Third defect on this page in the same family: the run before (11) also over-extended. Emphasis
% invented is as damaging as emphasis lost, and harder to notice, because a plausible pattern is
% exactly what a reader expects to see.
Die biquadratische Gleichung hat zwei Paare gleicher, reeller Wurzeln, wenn
\weq{8}{D = 0, \quad a < 0, \quad a^2-4c = 0}
und zwei Paare gleicher imaginärer Wurzeln, wenn
\weq{9}{D = 0, \quad a > 0, \quad a^2-4c = 0.}

Ist aber $\alpha$ von $\beta$ verschieden, so muss $\gamma$ reell sein und (7) zeigt, dass, wenn $\alpha$ und $\beta$ reell sind, $a$ negativ und $a^2-4c$ positiv ist, dass dagegen, wenn $\alpha$ und $\beta$ conjugirt imaginär sind, entweder $a$ positiv oder $a^2-4c$ negativ sein muss, da in diesem Falle $(\alpha-\beta)^2$ negativ ist, und wenn $2(\alpha+\beta)^2+(\alpha-\beta)^2$ positiv ist, jedenfalls auch $8(\alpha+\beta)^2+(\alpha-\beta)^2$ positiv sein muss.

Wir können also (6) dahin ergänzen, dass wir sagen:

\gesperrt{Die nothwendige und hinreichende Bedingung für die Existenz von vier reellen Wurzeln, von denen auch zwei (aber nicht mehr) einander gleich sein können, ist}
\weq{10}{D \geqq 0, \quad a < 0, \quad a^2-4c > 0.}

% [corrected 2026-08-07, 8x zoom strip 5/8] OVER-EXTENSION, the mirror of the usual defect: the
% wrap ran further than the print. The Sperrung stops after "für drei gleiche Wurzeln," --
% "die nothwendig alle reell sind, ist nach (2)" is set normally, and the contrast against the
% spaced words on the same line is plain at zoom.
%
% Note the statement immediately above (10) IS spaced in full, "ist" included. Two statements of
% the same grammatical shape, three lines apart, spaced to different extents: this is
% per-occurrence and must never be swept.
\gesperrt{Die nothwendige und hinreichende Bedingung für drei gleiche Wurzeln,} die nothwendig alle reell sind, ist nach (2)
\[
u^2 = v^2 = w^2,
\]
es muss also die cubische Resolvente (3) drei gleiche Wurzeln haben, und dafür sind nach \S.\,77 (1) die nothwendigen und hinreichenden Bedingungen
\weq{11}{\begin{aligned} a^2+12c &= 0\\ 2a^3-8ac+9b^2 &= 0. \end{aligned}}
Nach \S.\,47, (15), (16) sind die beiden Invarianten $A$ und $B$ der biquadratischen Gleichung
\[
A = a^2+12c, \quad B = 2a^3-72ac+27b^2.
\]
Also können wir die Gleichungen (11) auch so schreiben
\[
A = 0, \quad B+4aA = 0,
\]
so dass (11) gleichbedeutend ist mit
\[
A = 0, \quad B = 0.
\]
Durch Elimination von $c$ kann man aus (11) auch noch die Gleichung ableiten:
\weq{12}{8a^3+27b^2 = 0,}
